% !TeX spellcheck = it_IT
% !TeX root = ../scompl.tex

\subsection{Algoritmi Montecarlo e Las Vegas}

Definizioni di:
\begin{itemize}
	\item Algoritmo probabilistico
	
	\item Algoritmi Montecarlo: tempo deterministico, risultato possibilmente sbagliato
	\begin{itemize}
		\item One-sided: sempre giusto quando esce 1 ($q(I) = A(I,Z) = 1$)
		
		\item Two-sided: può sbagliare su entrambi gli output
	\end{itemize}
	
	\item Algoritmi Las Vegas: risultato corretto, tempo non deterministico
\end{itemize}

Amplificazione:
\begin{itemize}
	\item Montecarlo one sided: si può ridurre la probabilità di errore con esecuzioni indipendenti dell'algoritmo
	
	\item Montecarlo two-sided: voto di maggioranza su esecuzioni indipendenti, usando Chernoff-Hoeffding il voto di maggioranza è "probabilmente" corretto
\end{itemize}

Da Las Vegas a Montecarlo One-sided: 
\begin{itemize}
	\item Trovare $t(n) \geq c \cdot f(n)$, con $f(n)$ tempo di esecuzione massimo atteso dell'algoritmo Las Vegas
	
	\item Blocca l'esecuzione se l'algoritmo non termina in $t(n)$ passi, restituendo 0
	
	\item Sempre corretto se torna 1, errore con probabilità $\leq \frac{1}{c}$ per disuguaglianza di Markov
\end{itemize}

\subsection{Classi di complessità probabilistiche}

$\BPP$: problemi risolti efficientemente da algoritmi Montecarlo two-sided
$$ \Pr \left(B (I,Z) \neq q(I)\right) \leq \frac{1}{3} = \frac{1}{2} - \frac{1}{p'(|I|)} $$
\begin{itemize}
	\item $\P \subseteq \BPP$, $\P \stackrel{?}{\equiv} \BPP$
	
	\item $\BPP \stackrel{?}{\equiv} \NP$
	
	\item $\BPP \equiv \coBPP$: la definizione è simmetrica rispetto al complemento
\end{itemize}

$\RP$: problemi risolti efficientemente da algoritmi Montecarlo one-sided
$$
\begin{array}{r l}
	q(I) = 1 \implies & \displaystyle \Pr \left(B(I, Z) = 1\right) \geq \frac{2}{3} = \frac{1}{p'(|I|)} \\
	q(I) = 0 \implies & \displaystyle \Pr \left(B(I, Z) = 0 \right) = 1 
\end{array}
$$
\begin{itemize}
	\item $\P \subseteq \RP$
	
	\item $\RP \subseteq \NP$: riscrivendo la def di $\NP$
	
	\item $\coRP \subseteq \coNP$, $\P \subseteq \coRP$
	
	\item $\RP \subseteq \BPP$, $\coRP \subseteq \BPP$
\end{itemize}

$\ZPP$: def come $\ZPP \equiv \RP \cap \coRP$, problemi risolti in tempo atteso limitato da un polinomio nella lunghezza dell'istanza da algoritmi Las Vegas
$$
\begin{array}{r l}
	q(I) = 1 \implies & \displaystyle \Pr \left(B(I, Z) = 1\right) \geq \frac{2}{3} \wedge \Pr (B' (I, Z') = 1) = 1 \\
	q(I) = 0 \implies & \displaystyle \Pr \left(B'(I, Z') = 0\right) \geq \frac{2}{3} \wedge \Pr (B' (I, Z') = 0) = 1
\end{array}
$$
\begin{itemize}
	\item Se $X \in \ZPP$, alg $A$ che su $I$ esegue $B$ e $B'$ terminando quando $B(I, Z) = 1$ o $B'(I, Z') = 0$, la prob che ciò non accada è $1/3$; usando il $\Ex$ della Geometrica, $3/2 < 2$ ripetizioni è il numero atteso
	
	\item $\ZPP \subseteq \RP$: si può passare da LV a MC-OS
	
	\item $\P \subseteq \ZPP$: dato che $\P \subseteq \RP$ e $\P \subseteq \coRP$
\end{itemize}

\subsection{Estrattore di Von Neumann}

Si vuole trasformare una sorgente non perfettamente casuale in una completamente casuale. Lancia una moneta truccata finché non si ottengono due valori diversi di seguito, le due sequenze possibili sono i due possibili risultati, hanno la stessa prob.

Numero di lanci atteso prima di una sequenza utile: 
\begin{itemize}
	\item Definiamo $Z_k = 1 \Leftrightarrow$ se i lanci $2k$ e $2k-1$ sono una sequenza utile
	
	\item Per valore atteso della Geometrica trovo il numero medio di lanci
\end{itemize}

\subsection{Il problema del Coupon Collector}

Vogliamo osservare gli $n$ possibili valori distinti, ognuno ha prob $\frac{1}{n}$. 

Vogliamo trovare il numero minimo di realizzazioni per osservare ciascun $a_i$ almeno una volta:
\begin{itemize}
	\item $N_1, \dots, N_n$, con $N_i$ numero di estrazioni aggiuntive per ottenere l'$i$-esimo valore distinto; sono tutte Geometriche
	
	\item Data la probabilità di osservare un nuovo valore calcolo 
	$$ \Ex [N] = \sum_{i = 1}^n N_i = \dots = n \left(\ln n + \Theta (1)\right) $$
\end{itemize}

\subsection{Reservoir Sampling}

Vogliamo una struttura dati che contenga $k$ elementi e dato uno stream in ingresso per ogni $t \geq k$, ognuno dei primi $t$ elementi è presente nella struttura con prob $\frac{k}{t}$. Modello streaming e spazio $\Theta (k)$.

Definizione dell'algoritmo.

Dimostriamo che $\forall t \geq k$ vale $\Pr\left(x_i \in R_t \right) = \frac{k}{t}$, $\forall i \leq t$:
\begin{itemize}
	\item Per induzione, base con $t = k$
	
	\item Step: con $t \geq k$ dimostriamo $\Pr \left(x_i \in R_{t+1}\right) = \frac{k}{t+1}$, $\forall i \leq t+1$
	\begin{align*}
		\Pr \left(x_i \in \R_{t+1}\right) & = \Pr \left(x_i \in R_t\right) \Pr (x_i \in R_{t+1} \mid x_i \in R_t ) \\ 
		& = \Pr \left(x_i \in R_t\right) \left( 1 - \Pr (x_i \notin R_{t+1} \mid x_i \in R_t) \right) \\
		& = \frac{k}{t} \cdot \left(1 - \frac{k}{t + 1} \frac{1}{k} \right) = \frac{k}{t+1}
	\end{align*}
\end{itemize}

\subsection{Algoritmo di Karger per il taglio minimo}

Dato un multigrafo non orientato, vogliamo il taglio che corrisponde a una partizione ammissibile con il costo minimo.

Definizione di Karger-Base e Karger.

Vogliamo prima trovare la prob che Karger-Base torni il taglio minimo: 
\begin{itemize}
	\item Prob che l'$i$-esimo arco non sia nel taglio $\Gamma^\ast$
	$$ p_i = \Pr \left(A_i \mid A_1, \dots A_{i-1}\right) = \dots \geq \frac{|V| - i - 1}{|V| - i + 1} $$
	
	\item Prob che l'alg restituisca $\Gamma^\ast$
	$$ \Pr (\text{ret } \Gamma^\ast) = \prod_{i = 1}^{|V| - 2} p_i \geq \dots = \frac{1}{\binom{|V|}{2}} $$
\end{itemize}

Possiamo amplificare l'alg per ottenere probabilità arbitrariamente bassa:
\begin{itemize}
	\item Probabilità che fallisca su $M$ esecuzioni indipendenti
	$$ \Pr (\text{fail}) = \Pr (\text{ret } \Gamma^\ast)^M \leq e^{- \frac{1}{\binom{|V|}{2}} M} = \epsilon$$
	scegliendo $M = \left\lceil \binom{|V|}{2} \ln \frac{1}{\epsilon} \right\rceil$
\end{itemize}

Il tempo totale di esecuzione è $\O \left(|E| |V|^2 \ln \frac{1}{\epsilon}\right)$.